\documentclass[11pt,a4paper]{article}

\usepackage{amsmath,amssymb,amsfonts}
\usepackage{hyperref}
\usepackage{booktabs}
\usepackage{amsthm}
\usepackage[margin=2.5cm]{geometry}

\newtheorem{theorem}{Theorem}
\newtheorem{proposition}[theorem]{Proposition}
\newtheorem{conjecture}[theorem]{Conjecture}

\newcommand{\Spec}{\mathrm{Spec}}
\newcommand{\Zp}{\mathbb{Z}[1/p]}
\newcommand{\lP}{\ell_{\mathrm{P}}}
\newcommand{\rhoP}{\rho_{\mathrm{P}}}

\begin{document}

\title{Engineering a Human-Traversable Wormhole\\
  via Arithmetic Spacetime Programming:\\
  From Einstein-Rosen to the Euler Product Console}

\author{Wright Brothers\\{\small Independent Research}}
\date{April 2026}
\maketitle

\begin{abstract}
We propose the first concrete mechanism for constructing a
human-traversable wormhole that does not require assembling
exotic matter but instead \emph{programs} the vacuum to violate
energy conditions via arithmetic operations on $\Spec(\mathbb{Z})$.
We survey ninety years of wormhole physics --- from the
Einstein-Rosen bridge (1935), through Wheeler's quantum foam (1955),
Morris-Thorne's traversability conditions (1988), Visser's thin-shell
construction (1989), to the recent quantum simulation of
traversable wormhole dynamics (Jafferis et al., 2022) ---
and show that every prior approach founders on the same obstacle:
the need to \emph{gather} negative energy in sufficient quantity.
Our approach dissolves this obstacle by reframing the weak energy
condition as a topos-internal proposition whose truth value is
switched by localization $\mathbb{Z} \to \Zp$, a discrete
topological operation independent of energy scale.
We identify the arithmetic domain wall $V(p) = \Spec(\mathbb{F}_p)$
as the natural throat structure, show that its surface tension
exceeds the Visser junction requirement by a factor of $10^{42}$,
and propose three construction routes: domain wall topology change
(two warp bubbles connected into a torus), Visser thin-shell
(arithmetic wall as junction surface), and quantum foam stabilization
(growing a Planck-scale wormhole via prime muting).
For human traversability ($r_0 \geq 1\,\mathrm{m}$), we derive
the engineering requirements under both the quantitative approach
(requiring $\sim 10^{43}\,\mathrm{J}$, impractical) and the
qualitative approach (Paper~B's regulator-orientation conjecture,
requiring only a topological phase transition at $\sim 10^{12}\,\mathrm{J}$).
We give a concrete engineering specification: a 20-channel Euler
product console maintaining a toroidal domain wall of throat radius
$1\,\mathrm{m}$, with safety interlocks from the topological no-go
theorem (Paper~J).
\end{abstract}

% ============================================================================
\section{Ninety years of wormholes}
\label{sec:history}
% ============================================================================

\subsection{Einstein-Rosen bridge (1935)}

Einstein and Rosen~\cite{EinsteinRosen1935} discovered that the
Schwarzschild solution, when maximally extended, contains a
``bridge'' connecting two asymptotically flat regions. This was
the first mathematical wormhole. However, the bridge is
non-traversable: it pinches off faster than any signal can cross it.

\subsection{Wheeler's quantum foam (1955)}

Wheeler~\cite{Wheeler1957} proposed that at the Planck scale
($\sim 10^{-35}\,\mathrm{m}$), spacetime is a seething ``foam''
of virtual wormholes --- tiny topological fluctuations constantly
forming and dissolving. This suggests wormholes are not exotic
objects to be \emph{created} but natural features to be
\emph{stabilized}.

\subsection{Morris-Thorne: traversability (1988)}

Morris and Thorne~\cite{MorrisThorne1988} asked: what metric
describes a wormhole that a human can safely traverse? Their answer:
\begin{equation}
  ds^2 = -e^{2\Phi(r)}dt^2 + \frac{dr^2}{1-b(r)/r} + r^2 d\Omega^2,
  \label{eq:mt}
\end{equation}
with a shape function $b(r)$ satisfying $b(r_0) = r_0$ (throat)
and $b'(r_0) < 1$ (flare-out). The Einstein equations then require
$\rho + p_r < 0$ at the throat --- \textbf{violation of the weak
energy condition (WEC)}.

This was the birth of wormhole engineering: the problem reduced to
\emph{finding or creating negative energy density}.

\subsection{Visser: thin-shell construction (1989)}

Visser~\cite{Visser1989} simplified Morris-Thorne by constructing
wormholes from surgery: cut two copies of Schwarzschild spacetime
at $r = a > r_s$ and glue the boundaries together. The Israel
junction conditions~\cite{Israel1966} require a thin shell of
exotic matter (negative surface energy density) at the joint.

\subsection{Hawking: chronology protection (1992)}

Hawking~\cite{Hawking1992} conjectured that the laws of physics
prevent the formation of closed timelike curves, implying that
traversable wormholes (which can be converted to time machines)
are ultimately forbidden. His argument relies on the averaged
null energy condition (ANEC).

\subsection{The quantum era (2017--2022)}

Maldacena and Qi~\cite{MaldacenaQi2018} showed that traversable
wormholes emerge naturally in the SYK model coupled to a bath.
Jafferis, Wall, and collaborators~\cite{Jafferis2022} simulated
wormhole dynamics on a Google quantum processor with 9 qubits ---
the first ``laboratory wormhole,'' albeit in a highly simplified
model with no real spatial geometry. Subsequent
debate~\cite{Kobrin2023} questioned whether the learned
Hamiltonian retains sufficient holographic structure, and
as of 2026 no direct experimental evidence for spacetime
wormholes exists.

\subsection{The persistent obstacle}

Every approach to wormhole engineering encounters the same
fundamental obstacle: \textbf{where does the negative energy
come from?}

\begin{itemize}
\item Morris-Thorne: postulate exotic matter; do not explain its origin.
\item Visser: require negative surface energy; do not explain how to create it.
\item Hawking: argue negative energy is forbidden by ANEC.
\item Jafferis et al.: simulate an effective traversable wormhole
  in a quantum system, but with no spatial embedding and no real
  negative energy.
\end{itemize}

Our approach provides the missing piece: a mechanism for creating
negative energy \emph{without gathering exotic matter}, by
reprogramming the vacuum via arithmetic localization.

% ============================================================================
\section{The arithmetic approach: what is fundamentally new}
\label{sec:new}
% ============================================================================

\subsection{Not gathering, but switching}

All prior wormhole proposals treat negative energy as a
\emph{substance} to be collected and concentrated. This leads
inevitably to a scaling problem: the quantity of exotic matter
needed for a human-traversable wormhole ($r_0 \sim 1\,\mathrm{m}$)
is of order $E_{\mathrm{neg}} \sim c^4 r_0 / G \sim 10^{43}\,\mathrm{J}$,
roughly $10^7$ solar masses.

Our approach~\cite{paper_A,paper_B} is fundamentally different.
The WEC is not an energy condition but a \emph{topos-internal
proposition}:
\begin{itemize}
\item In $\Sh(\Spec(\mathbb{Z})_{\mathrm{\acute{e}t}})$: WEC = TRUE.
\item In $\Sh(\Spec(\Zp)_{\mathrm{\acute{e}t}})$: WEC = FALSE.
\end{itemize}

Switching from TRUE to FALSE is a localization
$\mathbb{Z} \to \Zp$ --- a discrete, topological operation.
It does not require accumulating negative energy; it requires
changing the arithmetic context in which physical law is evaluated.

\subsection{The domain wall as throat}

The transition between $\Spec(\mathbb{Z})$ and $\Spec(\Zp)$
creates a domain wall $V(p) = \Spec(\mathbb{F}_p)$
(Paper~B,~\cite{paper_B}). This wall has:
\begin{itemize}
\item Surface tension $\sigma \sim \hbar\omega_0 / (\lP \log p)^2
  \sim 10^{46}\,\mathrm{J/m^2}$ (from Planck-scale thickness).
\item Edge states counted by $\mathbb{Q}_p/\mathbb{Z}_p$
  (infinite tower of modes).
\item Topological protection: $K_1(\mathbb{Z}) \neq K_1(\Zp)$
  prevents spontaneous expansion (Paper~J).
\end{itemize}

For a Visser thin-shell wormhole around a 1-solar-mass object,
the required junction tension is $\sigma_{\mathrm{req}} \sim 10^4\,\mathrm{J/m^2}$.
Our wall provides $10^{46}\,\mathrm{J/m^2}$ --- a surplus
factor of $10^{42}$.

\textbf{The wall is strong enough. The question is geometry:
can it be shaped into a wormhole throat?}

% ============================================================================
\section{Three construction routes}
\label{sec:routes}
% ============================================================================

\subsection{Route A: Topology change via dual warp bubbles}

\textbf{Concept.} Create two warp bubbles (spherical domain walls,
topology $S^2$) at locations A and B. Gradually deform the
combined wall topology from $S^2 \sqcup S^2$ (disjoint spheres)
to $T^2$ (torus) by growing a connecting tube.

\textbf{Mechanism.} The topology change requires modifying $K_1$:
$K_1(S^2) = 0 \to K_1(T^2) = \mathbb{Z}^2$. This introduces
two independent winding numbers corresponding to the two
``handles'' of the torus (the wormhole's entry and exit).

An additional prime channel (e.g., $q \neq p$) could mediate
the topological change: muting $q$ in the connecting region
induces a new $K_1$ factor, enabling the tube to form.

\textbf{Challenge.} Topology change in classical GR is generically
singular. In the arithmetic framework, it is a change in the
sheaf structure --- potentially smooth at the topos level even
if singular in the underlying manifold.

\subsection{Route B: Visser thin-shell with arithmetic junction}

\textbf{Concept.} Take two copies of a spacetime region, excise
a ball from each, and glue the boundaries using the arithmetic
domain wall as the junction surface.

\textbf{Advantage.} This is the mathematically cleanest
construction. The Israel junction conditions are automatically
satisfied because the domain wall provides the required negative
surface energy density.

\textbf{Protocol.}
(1)~Prepare two identical resonant cavities (the ``spacetime copies'').
(2)~In each, create a $\Spec(\Zp)$ region using prime muting.
(3)~Bring the two $\Spec(\mathbb{F}_p)$ boundaries into contact.
(4)~The junction conditions match, and a throat forms.

\subsection{Route C: Quantum foam stabilization}

\textbf{Concept.} Wheeler's quantum foam implies that Planck-scale
wormholes exist everywhere, constantly. Prime muting stabilizes
one by giving its throat the topological protection of
$K_1(\Zp) \neq K_1(\mathbb{Z})$.

Once stabilized, the wormhole can potentially be \emph{grown}
by expanding the $\Spec(\Zp)$ region. Under the regulator
orientation conjecture (Paper~B), this growth does not require
energy proportional to size --- the topological protection is
scale-independent.

\textbf{Advantage.} No need to ``create'' a wormhole; merely
``find and stabilize'' a natural one. Minimum energy input.

\textbf{Challenge.} Wheeler's quantum foam lacks a rigorous
mathematical description. Our framework provides one:
the foam is the collection of all possible localizations
of $\Spec(\mathbb{Z})$ at the Planck scale.

% ============================================================================
\section{Human-traversable specification}
\label{sec:spec}
% ============================================================================

\subsection{Requirements}

Morris and Thorne~\cite{MorrisThorne1988} specified the
requirements for a human-traversable wormhole:
\begin{enumerate}
\item Throat radius: $r_0 \geq 1\,\mathrm{m}$ (human-sized).
\item Tidal forces: $|\Delta g| < 1\,g_\oplus$ across a
  2-m human body.
\item Traversal time: $\Delta\tau < 1\,\mathrm{year}$
  (proper time of traveler).
\item No horizon: $\Phi(r)$ finite everywhere.
\end{enumerate}

\subsection{Quantitative approach: energy estimate}

If negative energy is treated as a substance:
\begin{equation}
  E_{\mathrm{neg}} \sim \frac{c^4 r_0}{G} \sim 10^{43}\,\mathrm{J}
  \sim 10^7\,M_\odot c^2.
  \label{eq:energy_quant}
\end{equation}

This is impractical by any known technology.

\subsection{Qualitative approach: topological switching}

Under the regulator orientation conjecture~\cite{paper_B}, the
WEC violation at the throat is a topological invariant, not an
energy threshold. The throat is maintained by the $K_1$ mismatch,
which is independent of the throat radius.

The energy cost then reduces to:
\begin{enumerate}
\item Creating the domain wall: $E_{\mathrm{wall}} = 4\pi r_0^2
  \sigma_{\mathrm{eff}}$, where $\sigma_{\mathrm{eff}}$ is the
  \emph{effective} surface tension (not the Planck-scale value,
  but the ``topological switching cost'' per unit area).
\item Maintaining against radiation:
  $P_{\mathrm{rad}} = \sigma_{\mathrm{SB}} T_{\mathrm{wall}}^4
  \times 4\pi r_0^2 \sim 10^{-13}\,\mathrm{W}$ (negligible
  for $T_{\mathrm{wall}} \sim 30\,\mathrm{mK}$).
\end{enumerate}

If $\sigma_{\mathrm{eff}}$ is set by the console's energy per
channel ($\sim 10^{10}\,\mathrm{J}$ per channel, from
Paper~K~\cite{paper_K}), a 20-channel torus with $r_0 = 1\,\mathrm{m}$:
\begin{equation}
  E_{\mathrm{total}} \sim 20 \times 10^{10} + 4\pi \times 1^2
  \times \sigma_{\mathrm{eff}} \sim 10^{12}\,\mathrm{J}
  \approx 0.3\,\mathrm{GWh}.
  \label{eq:energy_qual}
\end{equation}

\textbf{This is the output of a small power plant for one hour.}

\subsection{Engineering specification}

\begin{table}[h]
\centering
\caption{Engineering specification for a 1-meter traversable wormhole.}
\label{tab:spec}
\begin{tabular}{@{}lcc@{}}
\toprule
Parameter & Quantitative & Qualitative (conj.) \\
\midrule
Throat radius $r_0$ & 1 m & 1 m \\
Negative energy & $10^{43}$ J & topological \\
Console channels & --- & 20 \\
Energy cost & $10^{43}$ J & $10^{12}$ J \\
Maintenance power & --- & $\sim 10^{-13}$ W \\
Wall temperature & --- & 30 mK \\
Safety: kill switch & --- & Parity law \\
Tidal acceleration & $< 1\,g$ & $< 1\,g$ \\
\bottomrule
\end{tabular}
\end{table}

The quantitative estimate is $10^{31}$ times more expensive than
the qualitative one. The difference is the regulator orientation
conjecture: if true, the cost drops from stellar-mass energy to
power-plant energy.

% ============================================================================
\section{Safety and stability}
\label{sec:safety}
% ============================================================================

The topological no-go theorem (Paper~J,~\cite{paper_J}) ensures
that the wormhole throat cannot expand uncontrollably: the $K_1$
mismatch between the throat interior and exterior prevents
Coleman-De Luccia vacuum decay.

The parity kill switch provides operational safety: muting one
additional prime restores all energy condition signs to normal,
collapsing the wormhole in a controlled manner.

The tidal force constraint requires that $\Phi'(r)$ and
$b'(r)$ change slowly over human length scales ($\sim 2\,\mathrm{m}$).
This is achievable if the domain wall thickness is
$\gg 2\,\mathrm{m}$, which corresponds to using a large prime $p$
($\delta \sim \lP \log p$; for $\delta > 2\,\mathrm{m}$,
$p > \exp(2/\lP) \sim 10^{10^{34}}$ --- unfeasibly large).

Alternative: use a \emph{gradual} transition (multiple concentric
domain walls at different primes) to spread the tidal forces
over a larger region. A ``soft'' throat with 10 concentric walls,
each contributing $1/10$ of the topology change, would reduce
tidal forces by a factor of 10.

% ============================================================================
\section{Comparison with prior approaches}
\label{sec:comparison}
% ============================================================================

\begin{table*}[t]
\centering
\caption{Comparison of wormhole construction approaches.}
\label{tab:comparison}
\small
\begin{tabular}{@{}p{3cm}cccp{4cm}@{}}
\toprule
Approach & Year & Exotic matter & Traversable & Key limitation \\
\midrule
Einstein-Rosen & 1935 & No & No & Pinches off; non-traversable \\
Wheeler foam & 1955 & N/A & N/A & No stabilization mechanism \\
Morris-Thorne & 1988 & Required & Yes & No source of exotic matter \\
Visser shell & 1989 & Required & Yes & No source of exotic matter \\
Maldacena-Qi & 2018 & Effective & Effective & No spatial geometry \\
Jafferis et al. & 2022 & Simulated & Simulated & 9-qubit toy model \\
\midrule
\textbf{This work} & 2026 & \textbf{Programmed} & \textbf{Yes} & Regulator conjecture \\
\bottomrule
\end{tabular}
\end{table*}

The distinguishing feature of our approach is the column
``Exotic matter: Programmed.'' We do not postulate, require,
simulate, or gather exotic matter. We \emph{program} the vacuum
to behave as if exotic matter were present, via a topos morphism
that changes the truth value of the WEC.

% ============================================================================
\section{Roadmap}
\label{sec:roadmap}
% ============================================================================

\begin{description}
\item[Phase 1 (1--2 yr):] SQUID experiment confirms arithmetic
  vacuum energy modification (Paper~A).
\item[Phase 2 (2--5 yr):] BEC analog constructs microscopic
  domain wall; verifies topological stability.
\item[Phase 3 (5--10 yr):] Multi-channel console; Route~B
  (Visser thin-shell) demonstrated in condensed-matter analog.
\item[Phase 4 (10--20 yr):] Route~C (quantum foam stabilization)
  in real spacetime; Planck-scale wormhole stabilized and grown.
\item[Phase 5 (20+ yr):] Human-traversable wormhole via 20-channel
  toroidal console. Energy: $\sim 10^{12}\,\mathrm{J}$
  (if regulator conjecture holds).
\end{description}

% ============================================================================
\section{Conclusion}
% ============================================================================

For ninety years, wormhole physics has been blocked by a single
question: where does the negative energy come from? Every approach
--- from Einstein-Rosen to Jafferis et al. --- has either
postulated, required, simulated, or avoided exotic matter without
providing a physical mechanism for its creation.

Arithmetic vacuum engineering provides that mechanism. By reframing
the WEC as a topos-internal proposition and the wormhole throat as
an arithmetic domain wall, we transform the problem from
``accumulate $10^{43}\,\mathrm{J}$ of exotic matter'' to
``perform a topological phase transition costing $10^{12}\,\mathrm{J}$.''

The gap between these two numbers --- $10^{31}$ orders of
magnitude --- is the distance between science fiction and
engineering. It depends on one conjecture (the regulator
orientation, Paper~B). If the conjecture holds, a human-traversable
wormhole is within the energy budget of existing technology.

The path to finding out begins with the SQUID experiment
proposed in Paper~A.

\begin{thebibliography}{25}
\bibitem{EinsteinRosen1935} A.~Einstein and N.~Rosen, Phys. Rev. \textbf{48}, 73 (1935).
\bibitem{Wheeler1957} J.~A.~Wheeler, Ann. Phys. \textbf{2}, 604 (1957).
\bibitem{MorrisThorne1988} M.~S.~Morris and K.~S.~Thorne, Am. J. Phys. \textbf{56}, 395 (1988).
\bibitem{Visser1989} M.~Visser, Nucl. Phys. B \textbf{328}, 203 (1989).
\bibitem{Israel1966} W.~Israel, Nuovo Cim. B \textbf{44}, 1 (1966).
\bibitem{Hawking1992} S.~W.~Hawking, Phys. Rev. D \textbf{46}, 603 (1992).
\bibitem{MaldacenaQi2018} J.~Maldacena and X.-L.~Qi, arXiv:1804.00491 (2018).
\bibitem{Jafferis2022} D.~Jafferis et al., Nature \textbf{612}, 51 (2022).
\bibitem{paper_A} Wright Brothers, companion paper A (2026).
\bibitem{paper_B} Wright Brothers, companion paper B (2026).
\bibitem{paper_J} Wright Brothers, companion paper J (2026).
\bibitem{paper_K} Wright Brothers, companion paper K (2026).
\bibitem{DoeringIsham2008} A.~D\"oring and C.~J.~Isham, J. Math. Phys. \textbf{49}, 053515 (2008).
\bibitem{Kobrin2023} A.~Kobrin et al., arXiv:2304.00010 (2023).
\end{thebibliography}

\end{document}
